\section*{Problem 3: Searching for Primitive Polynomials (Level 3)}

\subsubsection*{Mathematical part}
\textbf{Recall:} To construct an extension field $GF(p^n)$ with $p^n$ elements, we must define arithmetic modulo a polynomial of degree $n$. However, not just any polynomial will do. We need a \emph{primitive} polynomial $f(x)$.

If $f(x)$ is primitive, then the powers of $x$ modulo $f(x)$:
\[
1, x, x^2, x^3, \ldots
\]
will visit every non-zero element of the field exactly once before returning to 1. Since there are exactly $m = p^n - 1$ non-zero elements, we require:
\[
x^m \equiv 1 \pmod{f(x)}
\]
and absolutely no smaller positive power should equal one. 

Number theory provides a shortcut to test this:
\begin{enumerate}
    \item Verify that $x^m \equiv 1 \pmod{f(x)}$.
    \item Find all prime factors $q$ of $m$.
    \item Verify that $x^{m/q} \not\equiv 1 \pmod{f(x)}$ for all prime factors $q$.
\end{enumerate}
If both conditions hold, the polynomial is primitive.
\textbf{Note:} A primitive polynomial for $GF(p^n)$ \textbf{must} have degree $n$.


\subsubsection*{Implementation part}
\textbf{Task:} Implement \lstinline|is_primitive| using the polynomial arithmetic you built in Level 2.

\textbf{Details:} You are translating the mathematical test above directly into code. 
\begin{itemize}
    \item Use integer factorization to find all unique prime factors of $p^n - 1$.
    \item Construct $x$ as \lstinline|Polynomial([field.zero, field.one])| and $1$ as \lstinline|Polynomial([field.one])|.
    \item Use Python's built-in \lstinline|pow(base, exponent, modulus)|. The provided \lstinline|Polynomial.__pow__| boilerplate natively supports this optional third \lstinline|modulus| argument, enabling rapid modular binary exponentiation without evaluating massive intermediate polynomials.
\end{itemize}

\begin{center}
\renewcommand{\arraystretch}{1.5}
\begin{tabularx}{\textwidth}{@{} >{\bfseries}l X c @{}}
\toprule
Function & Arguments & Returns \\ \midrule
is\_primitive & \lstinline|poly: Polynomial| & \lstinline|bool| \\
\bottomrule
\end{tabularx}
\end{center}

\subsubsection*{Experimentation part}
\textbf{UI Guide:} In \lstinline|lab_dashboard.py|, you can inspect how different choices of polynomial affect the generated multiplication tables, and how a primitive polynomial perfectly cycles through all non-zero elements.

\vspace{1cm}
